\pdfbookmark[1]{II. Determination of the Nearest Integer}{complexe-II}
\begin{center}[II.]\par\medskip
DETERMINATION OF THE NEAREST INTEGER.
\end{center}
Let
\[
\begin{aligned}
\varepsilon^3&=1,\quad m=a+b\varepsilon+c\varepsilon\varepsilon\\
2a-b-c&=A+\alpha\\
2b-c-a&=B+\beta\\
2c-a-b&=C+\gamma
\end{aligned}
\]
where $A,B,C$ are integers and $\alpha,\beta,\gamma$ are positive proper fractions. We then have
\[
A+B+C+\alpha+\beta+\gamma=0
\]
\Needspace{5\baselineskip}
so there are three cases to distinguish:

I. $\alpha+\beta+\gamma=0$, hence $\alpha=0$, $\beta=0$, $\gamma=0$

1, $A\equiv B\equiv C\ (\text{mod. }3)$. Here $m$ itself is an integer.

2, $A-B\equiv B-C\equiv C-A\equiv\pm1\ (\text{mod. }3)$. Here $m\pm\frac{\varepsilon-\varepsilon\varepsilon}{3}\pdot\varepsilon^n$ is an integer.

II. $\alpha+\beta+\gamma=1$.

Here each of
\[
\begin{aligned}
&A+B\varepsilon+C\varepsilon\varepsilon+1\\
&A+B\varepsilon+C\varepsilon\varepsilon+\varepsilon\\
&A+B\varepsilon+C\varepsilon\varepsilon+\varepsilon\varepsilon
\end{aligned}
\]
\SourceLeaf{402}%
is divisible by $1-\varepsilon$, and one of these numbers by $3$. The quotient, or
\[
m+\frac{\varepsilon^n-\alpha-\beta\varepsilon-\gamma\varepsilon\varepsilon}{3}
\]
is the required integer.

III. $\alpha+\beta+\gamma=2$

Here
\[
\begin{aligned}
&A+B\varepsilon+C\varepsilon\varepsilon+\varepsilon+\varepsilon\varepsilon\\
&A+B\varepsilon+C\varepsilon\varepsilon+\varepsilon\varepsilon+1\\
&A+B\varepsilon+C\varepsilon\varepsilon+1+\varepsilon
\end{aligned}
\]
are divisible by $1-\varepsilon$, and one of these numbers by $3$. The quotient, or
\[
m+\frac{\varepsilon^n(\varepsilon+\varepsilon\varepsilon)-\alpha-\beta\varepsilon-\gamma\varepsilon\varepsilon}{3}
\]
is the required integer.

In all three cases the remainder has the form
\[
x+y\varepsilon+z\varepsilon\varepsilon
\]
so that $x,y,z$ are smaller than $\frac13$ in absolute value, and $x+y+z=0$. This necessarily gives
\[
xx+yy+zz=2xx-2yz=2yy-2xz=2zz-2xy<\tfrac29
\]
because two of the three quantities $x,y,z$ must have the same sign. Consequently, the determinant of the remainder is
\[
=\tfrac32(xx+yy+zz)<\tfrac13\qquad\text{Q. E. D.}
\]

The nearest integer is found \emph{more conveniently} as follows. Given $a+b\varepsilon+c\varepsilon\varepsilon=m$, set
\[
\begin{aligned}
b-a&=C+\gamma\\
c-b&=A+\alpha\\
a-c&=B+\beta
\end{aligned}
\]
where $A,B,C$ are the next smaller integers and $\alpha,\beta,\gamma$ are positive fractions.
There are three cases to distinguish:

\SourceLeaf{403}%
I. $\alpha+\beta+\gamma=0$; then $m$ itself is an integer

II. $\alpha+\beta+\gamma=1$; then the nearest integer is
\[
\begin{array}{ll}
B+(B+C)\varepsilon&\text{if }\alpha\text{ is the largest fraction.}\\
\qquad C\varepsilon+(A+C)\varepsilon\varepsilon&\qquad\beta\\
A+B\quad\pdot\quad+A\varepsilon\varepsilon&\qquad\gamma
\end{array}
\]

II. $\alpha+\beta+\gamma=2$; then the nearest integer is
\[
\begin{array}{ll}
B+1+(B+C+2)\varepsilon&\text{if }\alpha\text{ is the smallest fraction.}\\
\qquad(C+1)\varepsilon+(A+C+2)\varepsilon\varepsilon&\qquad\beta\\
A+B+2\quad\pdot\quad+(A+1)\varepsilon\varepsilon&\qquad\gamma
\end{array}
\]
In II, $1$, the remainder is $\beta+(\beta+\gamma)\varepsilon$, whose determinant is
\[
\begin{aligned}
&=\beta\beta+\beta\gamma+\gamma\gamma\\
&=\tfrac13-\tfrac16[(\alpha-\beta)(1+3\beta)+(\alpha-\gamma)(1+3\gamma)]
\end{aligned}
\]

\emph{More simply still:}

Arrange the fractions $a-[a],b-[b],c-[c]$ by size, and call them $p,q,r$ in that order. If all three are equal, $m$ is an integer. If they are unequal, let $t$ be any fraction between
\[
\begin{array}{llll}
p\text{ and }q,&\text{according as}&q-p&\text{is largest}\\
q\text{ and }r&&r-q&\\
r\text{ and }1+p&&1+p-r&
\end{array}
\]
Then
\[
[a-t]+[b-t]\varepsilon+[c-t]\varepsilon\varepsilon
\]
is the nearest integer.
\BookRule{.3\textwidth}

\SourceLeaf{404}%
\pdfbookmark[1]{III.}{complexe-III}
\begin{center}[III.]\end{center}
Let $\varepsilon^5=1$
\[
\begin{aligned}
a+b\varepsilon+c\varepsilon\varepsilon+d\varepsilon^2+e\varepsilon^4&=q'\\
a+b\varepsilon^{-1}+c\varepsilon^{-2}+d\varepsilon^{-3}+e\varepsilon^{-4}&=q''''
\end{aligned}
\]
\[
\begin{aligned}
(a-b)^2+(b-c)^2+(c-d)^2+(d-e)^2+(e-a)^2&=2p'\\
(a-c)^2+(b-d)^2+(c-e)^2+(d-a)^2+(e-b)^2&=2p''
\end{aligned}
\]
\[
\begin{aligned}
q'q''''&=-p'\varepsilon-p''\varepsilon\varepsilon-p''\varepsilon^3-p'\varepsilon^4=P'\\
q''q'''&=-p'\varepsilon\varepsilon-p''\varepsilon^4-p''\varepsilon-p'\varepsilon^3=P''
\end{aligned}
\]

\[
\text{Determinant}=P'P''=-p'p'+3p'p''-p''p''
\]
\[
\begin{aligned}
\text{Measure}&=2p'+2p''=2P'+2P''\\
&=5(aa+bb+cc+dd+ee)-(a+b+c+d+e)^2
\end{aligned}
\]
Multiplication by $1-\varepsilon$ gives the new measure $=8p'$
\[
\text{Maximum measure}=2\left(\frac{\sin72^\circ}{\sin36^\circ}+\frac{\sin36^\circ}{\sin72^\circ}\right)\surd D=4{,}472\surd D
\]

\emph{Modulus} $=1-\varepsilon$
\[
\begin{aligned}
1-\varepsilon&=x\\
\varepsilon&=1-x\\
\varepsilon\varepsilon&=1-2x+xx\\
\varepsilon^3&=1-3x+3xx-x^3\\
\varepsilon^4&=1-4x+6xx-4x^3+x^4\\
&=-4+6x-4xx+x^3
\end{aligned}
\]
Thus
\[
\begin{aligned}
\frac{1-\varepsilon^n}{1-\varepsilon}&\equiv n\text{ mod. }(1-\varepsilon)\\
\varepsilon^n&\equiv1-nx\text{ mod. }(1-\varepsilon)^2\\
\left(\frac{\varepsilon+\varepsilon^4}{\varepsilon\varepsilon+\varepsilon^3}\right)^n&\equiv1+nxx\text{ mod. }(1-\varepsilon)^3
\end{aligned}
\]
Thus a number that is $\equiv1\text{ mod. }(1-\varepsilon)^3$ can be a unit \emph{only} if it is also $\equiv1\ (\text{mod. }5)$.
\BookRule{.3\textwidth}

\SourceLeaf{405}%
\pdfbookmark[1]{IV. Some Remarks on the Measure of Numbers}{complexe-IV}
\begin{center}[IV.]\par\medskip
SOME REMARKS ON THE MEASURE OF NUMBERS.
\end{center}
Let $\varepsilon^n=1$, with $n$ prime,
\[
\begin{aligned}
m&=a+a'\varepsilon+a''\varepsilon\varepsilon+a'''\varepsilon^3+\ldots+a^{(n-1)}\varepsilon^{n-1}=f\varepsilon\\
D&=f\varepsilon\pdot f\varepsilon\varepsilon\pdot f\varepsilon^3\ldots f\varepsilon^{n-1}\\
f\varepsilon\pdot f\varepsilon^{n-1}&=-b'(\varepsilon+\varepsilon^{n-1})-b''(\varepsilon^2+\varepsilon^{-2})\\
&\qquad-b'''(\varepsilon^3+\varepsilon^{-3})\ldots
\end{aligned}
\]
then
\[
\begin{aligned}
2b'&=(a-a')^2+(a'-a'')^2+(a''-a''')^2+\text{etc.}\\
2b''&=(a-a'')^2+(a'-a''')^2+(a''-a'''')^2+\text{etc.}\\
&\text{etc.}
\end{aligned}
\]

Here $b',b'',b'''\dotstwo$ are therefore all positive quantities; they are called the \emph{partial measures} of $m$, and their sum
\[
\begin{aligned}
b'+b''+b'''+\text{etc.}={}&n(aa+a'a'+a''a''+\dotstwo)\\
&\qquad-(a+a'+a''+\text{etc.})^2
\end{aligned}
\]
the \emph{general measure}. If we set
\[
f\varepsilon\pdot f\varepsilon^{n-1}=c',\quad f\varepsilon\varepsilon\pdot f\varepsilon^{n-2}=c''\quad\text{etc.}
\]
then
\[
\begin{aligned}
&c'+c''+c'''+\text{etc.}+c^{\frac12(n-1)}\\
&\qquad=b'+b''+b'''+\text{etc.}+b^{\frac12(n-1)}
\end{aligned}
\]
\[
\begin{aligned}
&c'(\varepsilon+\varepsilon^{n-1})+c''(\varepsilon\varepsilon+\varepsilon^{n-2})+c'''(\varepsilon^3+\varepsilon^{n-3})+\text{etc.}\\
&\qquad=2(b'+b''+b'''+\text{etc.}+b^{\frac12(n-1)})-nb'
\end{aligned}
\]
\[
\begin{aligned}
&c'(2-\varepsilon-\varepsilon^{n-1})+c''(2-\varepsilon\varepsilon-\varepsilon^{n-2})\\
&\qquad+c'''(2-\varepsilon^3-\varepsilon^{n-3})+\text{etc.}=nb'
\end{aligned}
\]
\[
b'>\frac{n-1}{2n}(nD)^{\frac2{n-1}},\quad b'+b''+b'''+\text{etc.}>\frac{n-1}{2}\pdot D^{\frac2{n-1}}
\]

If, in general,
\[
f\varepsilon\pdot f\varepsilon^{n-1}=A+A'\varepsilon+A''\varepsilon\varepsilon+A'''\varepsilon^3+\dotstwo
\]
then the general measure is $\triangle=-A-A'-A''-\text{etc.}+nA$

Measure of $(1+\varepsilon)f\varepsilon\ldots$
\[
\triangle'=4\triangle-2n(A-A')=4\triangle-2nb'
\]
If $a+a'+a''+\dotstwo=0$, then $\triangle=n(aa+a'a'+a''a''+\text{etc.})$

and if $A+A'+A''+\ldots=0$, then $\triangle=nA$, $\triangle'=n(2A+2A')$

Thus, if one of the coefficients $A',A''$, etc., is negative and greater in absolute value than $\frac12A$, the measure can be reduced without changing the determinant.
\BookRule{.3\textwidth}

\SourceLeaf{406}%
\Needspace{13\baselineskip}
\pdfbookmark[1]{V.}{complexe-V}
\begin{center}[V.]\end{center}
If it should be confirmed that every unit is composed solely of factors of the form
\[
\frac{\varepsilon^\alpha-\varepsilon^\beta}{\varepsilon^\gamma-\varepsilon^\delta}
\]
the following theorem would be proved:

\emph{If $f(\varepsilon)$ is a unit, then}
\[
\frac{f(\varepsilon)}{f(\varepsilon^{-1})}=\varepsilon^n
\]

Even without \emph{assuming} \emph{that} theorem, the conclusion is easy to prove.
Let
\[
\frac{f\varepsilon}{f\varepsilon^{-1}}=F\varepsilon
\]
then
\[
F\varepsilon\pdot F\varepsilon^{-1}=1
\]
from which the theory of measure readily gives
\[
F\varepsilon=\pm\varepsilon^n
\]
The lower sign is impossible, however, since otherwise $f\varepsilon$ would have to be divisible by $1-\varepsilon$.

It is easy to prove that the determinant of a number different from $0$ cannot be $=0$. If the determinant is divisible by $m$, the number itself is divisible by $1-\varepsilon$. Consequently, if the determinant is divisible by $m^{m-1}$, the number itself must be divisible by $m$. This is absurd: with determinant $0$, the number could first be divided by $m$, without changing its determinant, until it was no longer divisible. The first assertion follows as follows. Let the given number be
\[
a+b\varepsilon+c\varepsilon\varepsilon+\text{etc.}\equiv a+b+c\dotstwo\text{ mod. }1-\varepsilon
\]
so its determinant is $\equiv(a+b+c\dotstwo)^{m-1}\text{ mod. }1-\varepsilon$.
\BookRule{.3\textwidth}

\SourceLeaf{407}%
\pdfbookmark[1]{VI.}{complexe-VI}
\begin{center}[VI.]\end{center}
Let $\varepsilon^n=1$
\[
\begin{aligned}
f\varepsilon&=a+b\varepsilon+c\varepsilon\varepsilon+d\varepsilon^3+\text{etc.}\\
m&=\text{determinant of this number}\\
\frac{m}{f\varepsilon}&=f\varepsilon\varepsilon\pdot f\varepsilon^3\ldots f\varepsilon^{n-1}\\
&=A+B\varepsilon+C\varepsilon\varepsilon+\text{etc.}=F\varepsilon
\end{aligned}
\]
The number $f\varepsilon$ corresponds to a root of the congruence $x^n\equiv1\ (\text{mod. }m.)$ Let it be $r$. We have
\[
\begin{aligned}
nA&=F1+F\varepsilon+F\varepsilon\varepsilon+\dotstwo\\
nB&=F1+\varepsilon^{-1}F\varepsilon+\varepsilon^{-2}F\varepsilon\varepsilon+\dotstwo\\
nC&=F1+\varepsilon^{-2}F\varepsilon+\varepsilon^{-4}F\varepsilon\varepsilon+\dotstwo\\
&\text{etc.}
\end{aligned}
\]

Thus, since $F\varepsilon\varepsilon,F\varepsilon^3,F\varepsilon^4$, etc., are divisible by $f\varepsilon$,
\[
\begin{gathered}
nA-F1-\varepsilon(nB-F1)\\
nA-F1-\varepsilon\varepsilon(nC-F1)\\
nA-F1-\varepsilon^3(nD-F1)\\
\text{etc.}
\end{gathered}
\]
are all divisible by $f\varepsilon$, or also
\[
\begin{gathered}
n(A-B)-\varepsilon n(B-C)\\
n(B-C)-\varepsilon n(C-D)\\
\text{etc.}
\end{gathered}
\]
are divisible by $f\varepsilon$; consequently [if $f\varepsilon$ is not divisible by $1-\varepsilon$, and $F\varepsilon$ is not divisible by a real integer]
\[
\varepsilon\equiv\frac{A-B}{B-C}\equiv\frac{B-C}{C-D}\equiv\frac{C-D}{D-E}\text{ etc. }(\text{mod. }f\varepsilon)
\]
\BookRule{.3\textwidth}

\clearpage
\SourceLeaf{408}%
\pdfbookmark[1]{Remarks on the Complex Numbers}{complexe-Bemerkungen}
\begin{center}REMARKS.\end{center}
\BookRule{.15\textwidth}
The investigations collected here under the common heading On the Theory of Complex Numbers are scattered notes in the manuscript. They contain the essential steps of the proof of \textsc{Fermat}'s theorem for the third and fifth powers. In incomplete notes that are not printed here, numbers composed of cube roots of unity are studied by means of both the theory of binary quadratic forms and the division of the circle. When applying the latter, during the preparation of the paper \emph{Disquisitionum circa aequationes puras ulterior evolutio}, the ternary cubic form into which $27\frac{x^n-1}{x-1}$ can be transformed for a prime $n\equiv1\text{ mod. }3$ was also set out, and the theory of composition of the related form $X^3+mY^3+mmZ^3-3mXYZ$ was developed at the same time.

The property of integers formed from cube roots of unity that is assumed in the investigations of Fragment [I], namely that each can be decomposed into prime factors in only one way, follows from \textsc{Euclid}'s procedure for finding the common divisors of two numbers, if one applies the theorem derived in [II] on the nearest integer to any given fraction.

\Needspace{11\baselineskip}
That the same fundamental property holds for numbers composed of fifth roots of unity follows from this: the fractional remainder formed by a rule entirely analogous to that in [II], either for $m$ or else for $m$ multiplied by a suitable unit $E$, is such that subtracting it from the given number $mE$ produces an integer, and its determinant does not exceed unity. As indicated in [III], the units can be derived from the theory of binary quadratic forms of determinant $5$, in conjunction with the decomposition of a real prime into four factors (for example, $11=\text{Det. }(2+\varepsilon)$), namely as products of powers of $\varepsilon$ and $1+\varepsilon$.
\begin{flushright}\textsc{Schering.}\end{flushright}
\BookRule{.3\textwidth}
